\subsection{DIP\_edge\_bnd\_porosity\_map.txt}\label{sec_Model!Hyperbolic!DIP!edgePorosityMap}
%\subfile{./FileFormatInputDIPCellPorosityMap}
\index{Model!Hyperbolic!DIP}
\subsubsection{Purpose}
The Dual Integral Porosity(DIP) model requires that the variations in the connectivity porosity $\phi_\Gamma$ be specified on an interface-per-interface basis.
%
\subsubsection{Format description}
The format is exactly the same as that of the cell porosity map, see \ref{sec_Model!Hyperbolic!DIP!cellPorosityMap}. The only difference is that the running index i goes from 1 to the number of boundary interfaces (not the number of cells).
\subsubsection{Sample files}
See Subsection \ref{sec_Model!Hyperbolic!DIP!cellPorosityMap}.
%%%
%%%
%%%
\subsection{DIP\_edge\_int\_porosity\_map.txt}\label{sec_Model!Hyperbolic!DIP!edgePorosityMap}
%\subfile{./FileFormatInputDIPCellPorosityMap}
\index{Model!Hyperbolic!DIP}
\subsubsection{Purpose}
The Dual Integral Porosity(DIP) model requires that the variations in the connectivity porosity $\phi_\Gamma$ be specified on an interface-per-interface basis.
%
\subsubsection{Format description}
The format is exactly the same as that of the cell porosity map, see \ref{sec_Model!Hyperbolic!DIP!cellPorosityMap}. The only difference is that the running index i goes from 1 to the number of internal interfaces (not the number of cells).
\subsubsection{Sample files}
See Subsection \ref{sec_Model!Hyperbolic!DIP!cellPorosityMap}.